\documentclass{article}
\usepackage[english]{babel}
\usepackage[utf8]{inputenc}
\usepackage{graphicx}

\title{Gamification in Education}
\author{Hassan Qadeer}
\date{\today}

\begin{document}
\maketitle

\section{Introduction}
\begin{quote}
    \itshape
    Once, we learned for knowledge.

    Then came the points.

    Then the badges.

    Then the leaderboards.

    And suddenly, education was no longer about learning.

    It was about XP.
\end{quote}

Gamification has become increasingly common in education, with points, badges, levels and rewards being used to motivate students and increase engagement. For Ester, however, gamification is not something she is particularly fond of.
Having worked in a school and as a private tutor, Ester has seen how deeply gamification can influence the learning experience. Her father is also a teacher and frequently talks about the use of different methods to engage students.

Today, Ester works as a consultant at a company that builds Salesforce CRM solutions. She has always been the kind of person who learns best on her own terms: kind and patient with others, but stubborn when someone tries to tell her how she ought to study. When her employer began pushing her to pass the Salesforce Administrator certification, she took the pressure seriously and put in the hours, but the badges, points and progress bars built into the learning platform never motivated her the way they were meant to. If anything, they made a subject she already respected feel like something being performed for a scoreboard rather than something she was learning for herself.
Gamification is therefore no longer limited to games. It has become part of education, professional development and everyday learning. This raises an important question: does gamification actually improve learning, or does it simply make education more engaging? 
Reviews of the field suggest that gamification does have positive cognitive, emotional and behavioural effects overall, but that the quality of the design and the specific game elements used matter as much as the presence of gamification itself \cite{rabah2018}.

\section{The Battle of the Gamers}
Ester is sitting at home, the room lit only by the pale glow of her laptop screen. 
On the desk beside her, a can of soda has gone flat, forgotten somewhere between the third and fourth round. 
The notification sound plays again, a short bright chime, the same chime that plays every time someone's name moves. 
She has learned to tell the difference between the chime that means someone climbed and the one that means someone fell, though the sound itself never changes, only her stomach does. 
On the screen, a leaderboard refreshes itself in real time. Names rearrange, numbers tick upward in green, small animated stars burst beside whoever sits at the top. 
Her own name sits near the bottom, grey against all that shine, and no matter how quickly she answers, no matter how many rounds she plays, it barely moves. 
Through the wall she can hear her brother in the next room, laughing, his keyboard clattering fast and easy, his name climbing without him even seeming to try. 
Downstairs, her mother calls up to ask how she is doing, and Ester types back that she is doing great, thumbs quick over the keys, the lie smoother than she expected it to be. Her palms are damp against the trackpad. 
Her jaw aches from smiling at a screen no one else can see. Every few minutes a new round begins, and the timer counts down in the corner, red numbers shrinking second by second, and her hand hovers over the keys, already bracing for the small drop her name will take before the numbers even land.

Somewhere close to midnight the chime rings one more time and Ester does not hear it. Her head has come to rest against her forearm, the cursor blinking uselessly in an empty answer field, and the room around her dissolves the way rooms do in sleep: quietly, without a single seam to mark where waking ends. She opens her eyes onto a stadium she has never been to, tiered seats rising into a sky the exact yellow of a Salesforce onboarding screen, and she knows, the way you sometimes know things in dreams without being told, that she is dreaming, and that this particular dream is not going to let her go until she has earned it.

Astro is waiting for her at the centre of the field. Not the small plush toy from the trade show booths, but life sized, a raccoon zipped inside its own furred hood, round black eyes catching no light at all. Its voice, when it speaks, is warm in the practiced way customer service is warm. Three towers rise out of the grass at its gesture, each one a level Ester will have to clear before she can go home. Argue well, Astro tells her, and she moves on. Fail to convince it, and she stays.

The first tower is Motivation. Astro goes first, almost gently: points and badges make people want to keep going, that is simply how they are built, engagement climbing with every reward. Ester answers with the thing she has been carrying all night without naming it until now, that wanting to keep clicking has never once felt the same as understanding, that a badge can measure how long she stayed but never what she actually learned while she was there. The tower dims by one level. She climbs.

The second tower is Control, and this one Astro argues harder for, because this one it genuinely believes: a well designed system nudges a learner toward the behaviour that helps them most, quickly and painlessly, and where is the harm in that. Ester thinks of every evening she has spent doing exactly what a progress bar told her to do instead of what she actually wanted to understand, and she gives Astro the word it clearly was not built to hear back: manipulation, the quiet but important difference between guiding someone and deciding for them, the same difference ethicists have been trying to draw a line around since gamification first left the arcade \cite{kim2016}. The tower dims. She climbs again.

The third tower is the one she already knows, because she has been living inside it all evening: the leaderboard. Astro calls it healthy competition. Ester calls it by what it actually does: it publishes every failure in real time and hands it a number, so that the most cited harm of gamified learning has never really been the losing itself, but being seen losing \cite{toda2018}. For a moment Astro says nothing at all.

The third tower goes dark, and so does the stadium, and so does Astro, folding inward like a costume with no one left inside it.

Ester wakes up at her desk. The laptop screen has dimmed itself to save power, and in its reflection she can just make out her own face, tired, unremarkable, hers. Downstairs, her mother is still awake. Ester closes the laptop before the leaderboard gets the chance to load again.

\nocite{*}
\bibliographystyle{ieeetr}
\bibliography{references}

\end{document}